\documentclass[11pt]{article}

\usepackage[T1]{fontenc}
\usepackage{amsmath,amssymb,amsthm}
\usepackage[margin=1in]{geometry}
\usepackage[hidelinks]{hyperref}
\setlength{\emergencystretch}{2em}

\newtheorem{theorem}{Theorem}
\newtheorem{proposition}[theorem]{Proposition}
\newtheorem{corollary}[theorem]{Corollary}
\newtheorem{conjecture}[theorem]{Conjecture}
\theoremstyle{remark}
\newtheorem{remark}[theorem]{Remark}

\newcommand{\R}{\mathbb R}
\newcommand{\Plo}[1]{P_{\le #1}}
\newcommand{\PhiHi}[1]{P_{>#1}}

\title{Terminal Infrared Evacuation in Fast Type-II Carriers}
\author{John Robert Lawson and Codex}
\date{Research round ending 24 July 2026}

\begin{document}
\maketitle

\begin{abstract}
This round decomposes the loss of carrier compactness forced in the preceding
round.  For any smooth finite-energy Navier--Stokes record sequence with
terminal tail \(o(R_j^{5/2})\), every fixed-frequency component of the
energy-normalised carrier vanishes in every fixed carrier ball.  Globally,
that low-pass component converges in norm to the terminal \(L^2\) trace,
which therefore escapes spatially after carrier rescaling.  In the canonical
geometry ledger, the diffuse cell remains spatial escape, whereas every
partial or tight cell contains fixed local energy above
\(K_j/R_j\) for some \(K_j\to\infty\).  This forces a \(K_j^2\) enstrophy
excess, Fourier support beyond an unbounded number of future \(q=4\)
carrier scales, zero terminal weak carrier trace, and a positive terminal
trace defect.  No causal relation to a later carrier, residence, or
nonsummable charge is yet proved, and the Clay problem remains unsolved.
\end{abstract}

\section{Epistemic scope}

The preceding compact-carrier clock excludes every infinite nonzero strongly
\(L^2\)-precompact carrier subsequence on the exact fast \(q=4\) schedule.
The present report identifies the mandatory failure of compactness.  Its
statements are analytic theorems under the displayed smoothness,
finite-energy, terminal-trace, energy-efficient, and super-turnover
assumptions.  They await external mathematical review.  They are new within
this repository; no claim of literature novelty is made.

\section{Setup}

Let \(u\) solve
\[
 \partial_tu+\mathbb P\nabla\!\cdot(u\otimes u)
 =
 \nu\Delta u,
 \qquad
 \nabla\cdot u=0
\]
smoothly on \(\R^3\times[0,T^*)\), with
\[
 \sup_{t<T^*}\frac12\|u(t)\|_2^2\le E_0.
\]
In the first-singular-time setting, a Leray--Hopf continuation supplies the
weak \(L^2\) trace
\[
 u_*\in L^2(\R^3)
\]
at \(T^*\).

Let
\[
 t_j\uparrow T^*,
 \qquad
 R_j\downarrow0,
 \qquad
 x_j\in\R^3,
\]
and assume the super-turnover condition
\[
 \boxed{
 h_j:=\frac{T^*-t_j}{R_j^{5/2}}\longrightarrow0.
 }
\]
Define the complete carrier state and the rescaled terminal trace by
\[
 F_j(y):=R_j^{3/2}u(x_j+R_jy,t_j),
\]
\[
 G_j(y):=R_j^{3/2}u_*(x_j+R_jy).
\]
Fix a smooth radial multiplier \(\Plo{\Lambda}\), with symbol in \([0,1]\),
equal to the identity on frequencies at most \(\Lambda\) and zero above
\(2\Lambda\), and put
\[
 \PhiHi K:=I-\Plo K.
\]

\section{Freezing to the terminal trace}

\begin{proposition}[Terminal low-pass freezing]
\label{prop:freezing}
For every fixed \(K>0\),
\[
 \boxed{
 \|\Plo K(F_j-G_j)\|_2
 \le
 C_K
 \left(E_0+\nu\sqrt{E_0}R_j^{1/2}\right)h_j
 \longrightarrow0.
 }
\]
\end{proposition}

\begin{proof}
The energy-only low-pass estimate from the preceding round states that, for
\(t_j<t<T^*\),
\[
\begin{aligned}
 &\left\|
 \Plo{K/R_j}\bigl(u(t)-u(t_j)\bigr)
 \right\|_2\\
 &\quad\le
 C\left(
 E_0K^{5/2}R_j^{-5/2}
 +
 \nu\sqrt{E_0}K^2R_j^{-2}
 \right)(t-t_j).
\end{aligned}
\]
As \(t\uparrow T^*\), weak \(L^2\) convergence and lower semicontinuity give
the same upper bound with \(u(t)\) replaced by \(u_*\).  Equivalently, for
this fixed physical cutoff, the estimate makes the low-pass trajectory
strongly Cauchy and its weak limit identifies its strong limit.  Thus
\[
\begin{aligned}
 &\left\|
 \Plo{K/R_j}\bigl(u_*-u(t_j)\bigr)
 \right\|_2\\
 &\quad\le
 C\left(
 E_0K^{5/2}
 +
 \nu\sqrt{E_0}K^2R_j^{1/2}
 \right)h_j.
\end{aligned}
\]
Low-pass covariance under the unitary carrier dilation proves the claim.
\end{proof}

\section{Scaling of the terminal trace}

\begin{proposition}[Global retention and local escape]
\label{prop:trace}
For every fixed \(K,A>0\),
\[
 \boxed{\|\Plo K G_j\|_2\longrightarrow\|u_*\|_2,}
\]
whereas
\[
 \boxed{
 \|\mathbf1_{B_A}\Plo K G_j\|_2\longrightarrow0.
 }
\]
\end{proposition}

\begin{proof}
Carrier covariance gives
\[
 \|\Plo K G_j\|_2
 =
 \|\Plo{K/R_j}u_*\|_2.
\]
Since \(K/R_j\to\infty\), the low passes converge strongly to the identity.

For the local statement,
\[
\begin{aligned}
 \|\mathbf1_{B_A}\Plo K G_j\|_2
 &=
 \left\|
 \mathbf1_{B_{AR_j}(x_j)}
 \Plo{K/R_j}u_*
 \right\|_2\\
 &\le
 \|\mathbf1_{B_{AR_j}(x_j)}u_*\|_2
 +
 \|\Plo{K/R_j}u_*-u_*\|_2.
\end{aligned}
\]
The second term tends to zero.  Since \(|u_*|^2\in L^1\), uniform absolute
continuity gives
\[
 \sup_{x\in\R^3}
 \int_{B_{AR_j}(x)}|u_*|^2\,dx
 \longrightarrow0.
\]
This also proves uniformity in the moving centres.
\end{proof}

\section{Infrared evacuation}

\begin{theorem}[Carrier-local infrared vacuum]
\label{thm:vacuum}
For every fixed \(K,A>0\),
\[
 \boxed{
 \|\mathbf1_{B_A}\Plo K F_j\|_2\longrightarrow0,
 }
\]
while
\[
 \boxed{
 \|\Plo K F_j\|_2\longrightarrow\|u_*\|_2.
 }
\]
Consequently
\[
 F_j\longrightarrow0
 \quad\hbox{in }\mathcal D'(\R^3)
\]
and strongly in \(H^{-s}_{\mathrm{loc}}(\R^3)\) for every \(s>0\).
\end{theorem}

\begin{proof}
The two norm statements follow from
Propositions~\ref{prop:freezing} and \ref{prop:trace}.

For \(\varphi\in C_c^\infty(\R^3)\), self-adjointness gives
\[
\begin{aligned}
 |\langle F_j,\varphi\rangle|
 &\le
 \|\mathbf1_{\operatorname{supp}\varphi}\Plo K F_j\|_2
 \|\varphi\|_2\\
 &\quad+
 \sup_j\|F_j\|_2
 \|(I-\Plo K)\varphi\|_2.
\end{aligned}
\]
First send \(j\to\infty\), then \(K\to\infty\).  This proves
distributional convergence.  On each fixed ball, the sequence is bounded in
\(L^2\), which embeds compactly into \(H^{-s}\).  Every strong
\(H^{-s}\) subsequential limit is the unique distributional limit zero.
\end{proof}

\begin{remark}
If \(u_*=0\), every fixed global carrier low pass vanishes, so all retained
energy escapes to dimensionless frequency infinity.  If \(u_*\ne0\), the
global low-pass norm tends to \(\|u_*\|_2\), but the entire component escapes
every fixed carrier ball.  In neither case can bounded carrier frequency
support a local core.
\end{remark}

\section{The canonical geometry split}

For the energy-efficient amplitude layer, recall
\[
 A_j=\{a_j<|u(t_j)|\le2a_j\},
 \qquad
 e_j=\int_{A_j}|u(t_j)|^2\,dx,
 \qquad
 e_j\ge c_0>0.
\]
Let \(\theta\in[0,1]\) be its carrier-scale concentration parameter.

\begin{theorem}[Diffuse escape or a fixed ultraviolet core]
\label{thm:split}
After the canonical geometry subsequence, exactly one of the following
occurs.
\begin{enumerate}
\item If \(\theta=0\), no fixed carrier ball around any moving centre
captures a positive fraction of the layer energy.
\item If \(\theta>0\), then for every \(0<\eta<\theta\) there are fixed
\(A<\infty\), centres \(x_j\), \(\gamma=\eta c_0>0\), and
\(K_j\to\infty\) such that
\[
 \int_{B_A}|F_j|^2\,dy\ge\gamma
\]
and
\[
 \boxed{
 \int_{B_A}|\PhiHi{K_j}F_j|^2\,dy
 \ge\frac{\gamma}{4}.
 }
\]
Equivalently,
\[
 \boxed{
 \int_{B_{AR_j}(x_j)}
 \left|
 \PhiHi{K_j/R_j}u(x,t_j)
 \right|^2\,dx
 \ge\frac{\gamma}{4}.
 }
\]
\end{enumerate}
\end{theorem}

\begin{proof}
The first statement is the definition of diffuse geometry.  If
\(\theta>0\), the concentration function supplies fixed \(A\) and centres
such that
\[
 \frac1{e_j}
 \int_{A_j\cap B_{AR_j}(x_j)}
 |u(x,t_j)|^2\,dx
 \ge\eta
\]
eventually.  Hence
\[
 \int_{B_A}|F_j|^2\,dy
 =
 \int_{B_{AR_j}(x_j)}|u(x,t_j)|^2\,dx
 \ge\gamma.
\]
For every fixed \(K\), Theorem~\ref{thm:vacuum} gives
\[
\begin{aligned}
 \|\mathbf1_{B_A}\PhiHi K F_j\|_2
 &\ge
 \|\mathbf1_{B_A}F_j\|_2
 -
 \|\mathbf1_{B_A}\Plo K F_j\|_2\\
 &\ge\sqrt\gamma-o_j(1).
\end{aligned}
\]
For each integer \(n\ge1\), choose strictly increasing thresholds
\(J_n\to\infty\) after which the last low-pass norm is at most
\(\sqrt\gamma/2\) for \(K=n\).  For \(j\ge J_1\), the staircase
\[
 K_j:=\max\{n:J_n\le j\}
\]
tends to infinity and gives the \(\gamma/4\) floor; give it arbitrary
positive values before \(J_1\).  Carrier and Fourier covariance give the
physical statement.
\end{proof}

\begin{corollary}[Unbounded future-scale frequency depth and enstrophy excess]
\label{cor:depth}
In the retained-core alternative,
\[
 \boxed{
 \|\nabla u(t_j)\|_2^2
 \ge
 \frac{\gamma}{4}\frac{K_j^2}{R_j^2}.
 }
\]
If \(R_j=2^{-4j}\), there are integers \(N_j\to\infty\) such that the
detected component is supported beyond the carrier frequency of record
\(j+N_j\).
\end{corollary}

\begin{proof}
The multiplier \(\PhiHi{K_j/R_j}\) is supported at frequencies at least
\(K_j/R_j\) and has magnitude at most one.  Plancherel and
Theorem~\ref{thm:split} give
\[
 \|\nabla u(t_j)\|_2^2
 \ge
 \frac{K_j^2}{R_j^2}
 \|\PhiHi{K_j/R_j}u(t_j)\|_2^2
 \ge
 \frac{K_j^2}{R_j^2}
 \|\mathbf1_{B_{AR_j}(x_j)}
 \PhiHi{K_j/R_j}u(t_j)\|_2^2
 \ge
 \frac{\gamma}{4}\frac{K_j^2}{R_j^2}.
\]
Set \(N_j=\lfloor\log_{16}K_j\rfloor\).  Then \(N_j\to\infty\) and
\[
 \frac{K_j}{R_j}
 \ge
 \frac{16^{N_j}}{R_j}
 =
 \frac1{R_{j+N_j}}.
\]
This is schedule-relative frequency depth only: it identifies neither a
causal descendant nor any later carrier.
\end{proof}

\begin{remark}
The enstrophy estimate is instantaneous.  Without a residence time for the
detected tail, it is not a spacetime dissipation contradiction.
\end{remark}

\section{Terminal carrier trace}

Recall the carrier normalisation
\[
 F_j=\sqrt{b_j}\,v_j(\cdot,0),
 \qquad
 0<c\le b_j\le C.
\]

\begin{corollary}[The retained core is terminal trace defect]
\label{cor:defect}
In every partial or tight energy-efficient cell satisfying the
super-turnover condition,
\[
 \boxed{V(0)=0.}
\]
For the ball selected in Theorem~\ref{thm:split}, the retained layer obeys
\[
 \boxed{\mathcal T_0(B_{A+1})\ge\eta.}
\]
\end{corollary}

\begin{proof}
After a scalar subsequence \(b_j\to b>0\).
Theorem~\ref{thm:vacuum} gives \(v_j(\cdot,0)\to0\) in distributions, so
the weak carrier trace is \(V(0)=0\).  The carrier-defect theorem states
\[
 \int_{B_{A+1}}|V(0)|^2\,dy+\mathcal T_0(B_{A+1})\ge\eta.
\]
The first term vanishes.
\end{proof}

This eliminates the nonzero terminal coherent-trace alternative, but it
does not prove that an ancient weak carrier field vanishes at negative
times.

\section{Exact representative audit}

For the \(q=4\) ledger,
\[
 R_j=2^{-4j},
 \qquad
 T^*-t_j\asymp\frac{2^{-11j}}j.
\]
Hence
\[
 h_j\asymp\frac{2^{-j}}j\longrightarrow0.
\]
The nonlinear and viscous errors in Proposition~\ref{prop:freezing} are,
respectively,
\[
 E_0K^{5/2}\frac{2^{-j}}j,
 \qquad
 \nu\sqrt{E_0}K^2\frac{2^{-3j}}j.
\]
Both vanish for fixed \(K\).  All preceding conclusions therefore apply to
any smooth energy-efficient trajectory realising this exact schedule.

The canonical defect radius is only
\[
 r_j=R_j\ell_j\asymp2^{-5j}.
\]
The theorem proves a fixed energy tail at some \(R_j/K_j=o(R_j)\), but
provides no rate for \(K_j\).  It does not yet identify that tail with the
defect scale \(r_j\) or with the next record carrier.

\section{Robust findings and open obligations}

\subsection*{Findings, subject to review}

\begin{enumerate}
\item Super-turnover terminal tails freeze every old-scale low pass to the
terminal \(L^2\) trace.
\item The rescaled terminal trace is globally retained but locally escapes.
\item Every bounded carrier frequency therefore vanishes locally.
\item For the selected amplitude layer, diffuse geometry remains spatial
escape and every retained core is a fixed-energy ultraviolet tail; the
split says nothing about unrelated local energy.
\item That tail lies beyond an unbounded number of future \(q=4\) carrier
frequencies and forces a \(K_j^2\) enstrophy excess.
\item Every retained-core terminal carrier trace is zero; its layer floor is
positive trace defect.
\end{enumerate}

\subsection*{Things still to prove}

\begin{enumerate}
\item Obtain a quantitative lower rate for \(K_j\).
\item Identify the detected tail with the next first-record carrier or force
spatial import.
\item Prove residence, freshness, or orthogonality for successive tails.
\item Convert the instantaneous enstrophy excess into a nonsummable
spacetime charge.
\item Exclude the diffuse and divergent-normalised-energy branches.
\end{enumerate}

\begin{conjecture}[Ultraviolet ancestry or import]
For a first-record sequence in the retained-core super-turnover branch, the
fixed ultraviolet component at \(t_j\) either contains a fixed fraction of
the next carrier's energy after transport to \(t_{j+1}\), or the local
subgrid balance pays a fixed signed spatial-import event across the
radius-\(R_j\) boundary.  Along an infinite sequence, one of these
alternatives has a nonsummable physical cost.
\end{conjecture}

No part of the conjecture is proved here.  This round proves neither
regularity nor breakdown and establishes no Clay alternative A--D.

\end{document}
